\bookStart{Spae of Griper}[Grípis spǫ́]
\setBookCode{Gripisspa}

\begin{flushright}%
\textbf{Dating} \parencite{Sapp2022}: early C11th (0.616)--late C11th (0.313).

\textbf{Meter:} \Fornyrdislag%
\end{flushright}

\section{Introduction}

The \textbf{Spae of Griper} (abbrev. \Gripisspa) is a long poem dealing with the life of Siward.  It serves as a summary of the majority of the Walsing-cycle.

\subsection{Textual preservation and transmission}

\Gripisspa\ is only preserved in \Regius, where it is found on foll.~26v/19--28v/29, following \textlink{HelgakvidaOne} and foregoing \textlink{Reginsmal}.

\subsection{Meter}

The meter is strict \Fornyrdislag.  Anacrusis is occasionally allowed (3/2b, 45/3a).  No half-lines of type A3 appear in the b-verse.
The length of the stanza is unusually regular; indeed, every single one of the 53 \Fornyrdislag\ stanzas is exactly 4 lines long, something otherwise attained only by the much shorter \textlink{GudrunThree}.

An unusual metrical trait of \Gripisspa\ is the occurrence of catalectic lines, likely due to the loss of hiatus.  After this had occurred, lines like TODO (< TODO: uncontracted form) and TODO (< TODO: form which was always monosyllabic) would have scanned identically, leading to the acceptance of the new, catalectic type as a genuine metrical type.  TODO: Cite Males on Hiatus.

\subsection{Dating}

Although \textcite{Sapp2022}’s model arrives at a most probable early 11th c.~date for the poem, this result contradicts with reliable linguistic diagnostic criteria.

TODO: Cite more examples and papers below.

Numerous traits speak to a late date for \Gripisspa.  Linguistically, the poem displays the sound change \emph{vr-} > \emph{r-} (26/1, 49/1), which indicates a date of composition after 1000 \ce.  It also shows hiatus contraction, which is not found in Scaldic poetry older than ca.~1150 \ce.  It has four metrically required occurrences of the short form \emph{Sig·urðr} (16/3, 24/2, 40/2, 53/4), a younger form which is not found in any other Eddic poem apart from \textlink{Hyndluljod}, another likely late composition \parencite[138--139]{Sievers1889}.  In addition, it uses the word \emph{ęigi} ‘not’, originally ‘not at all, by no means’, as an unmarked negator.

Apart from these signs, \Gripisspa\ displays influence from numerous poems in the Walsing-cycle (e.g., 17/4 n.), and serves as an overview for the entire story, which points towards a systematising tendency, something associated with antiquarians of the late 12th and 13th centuries.  The poem is also exceptionally regular with regards to the length of the stanza, which probably suggests that it has not been transmitted for very long before ending up in the archetype to the heroic half of \Regius.

Together, these traits point convincingly towards a composition in the late 12th, or even early 13th c.

\newpage

\section{Text}

\subsection{Frá dauða Sin·fjǫtla (‘From the Death of Sinfittle’)}

\sectionline

\bpg\bpa Sig·mundr Vǫlsungs sonr var konungr á Frakk-landi.  Sin·fjǫtli var eldstr hans sona, annarr Helgi, þriði \edtext{Há·mund*r}{\Afootnote{\emph{Hámundir} \Regius}}.  Borg·hildr, kona Sig·mundar, átti bróður er hét \edtext{...}{\Afootnote{An unfilled gap for a lost name is clearly indicated in \Regius}}  En Sin·fjǫtli, stjúp-sonr hennar, ok \edtext{...}{\Afootnote{A gap is indicated in \Regius}} báðu einnar konu báðir ok fyr þá sǫk drap Sin·fjǫtli hann.  En er hann kom heim þá bað Borg·hildr hann fara á brot en Sig·mundr bauð henni fé-bǿtr ok þat varð hón at þiggja.  En at erfi’nu bar Borg·hildr ǫl.  Hón tók eitr mikit, horn fullt, ok bar Sin·fjǫtla.  En er hann sá í horn’it skilði hann at eitr var í ok mę́lti til Sig·mundar: „Gjǫr-óttr er drykkr’inn, ái!“  Sig·mundr tók horn’it ok drakk af.  Svá er sagt at Sig·mundr var harð-gǫrr at hvár-ki mátti hánum eitr granda útan né innan.  En allir synir hans stóðu⸗sk eitr á hǫrund útan.  Borg·hildr bar annat horn Sin·fjǫtla ok bað drekka ok fór allt sem fyrr.  Ok enn it þriðja sinn bar hón hánum horn’it ok þó á·mę́lis-orð með ef hann drykki eigi af.  Hann mę́lti enn sem fyrr við Sig·mund; hann sagði: „Lát-tu grǫn sía þá, sonr!“  Sin·fjǫtli drakk ok varð þegar dauðr.
Sig·mundr bar hann langar leiðir í fangi sér ok kom at firði einum mjóvum ok lǫngum ok var þar skip eitt lítit ok maðr einn á.  Hann bauð Sig·mundi far of fjǫrð’inn.  En er Sig·mundr bar lík’it út á skip’it þá var bátr’inn hlaðinn.  Karl mę́lti at Sig·mundr skyldi fara fyr inn á fjǫrð’inn.  Karl hratt út skip’inu ok hvarf þegar.
Sig·mundr konungr dvalði⸗sk lengi í Danmǫrk í ríki Borg·hildar síðan er hann fekk hennar.  Fór Sig·mundr þá suðr í Frakk-land til þess ríkis er hann átti þar.  Þá fekk hann Hjǫr·dísar, dóttur Ey·lima konungs.  Þeira sonr var Sig·urðr.  Sig·mundr konungr fell í orrustu fyr Hundings sonum.  En Hjǫr·dís gipti⸗sk þá Álfi, syni Hjalp·reks konungs.  Óx Sig·urðr þar upp í barn-ǿsku.  Sig·mundr ok allir synir hans vóru langt um fram alla menn aðra um afl ok vǫxt ok hug ok alla at·gørvi.  Sig·urðr var þá allra framarstr ok hann kalla allir menn í forn-frǿðum um alla menn fram ok gǫfgastan her-konunga.\epa

\bpb Syemund, Walsing’s son, was a king in Frankland.  Sinfittle was the eldest of his sons, the other Hallow, the third Highmund.  Burhild, Syemund’s wife, had a brother whose name was ...  But Sinfittle, her stepson, and ... both asked for the hand of a single woman, and for that reason Sinfittle slew him.  But when he came home Burhild bade him go away, but Syemund offered her restitutory money and she had to accept it.  But at the wake Burhild brought forth the ale.  She took much venom, a full horn, and brought it to Sinfittle.  But when he looked into the horn he understood that there was venom in it, and said to Syemund: “Treacherous is this drink, father!”  Syemund took the horn and drank it up.  So it is said that Syemund was sturdy enough that venom could not injure him from without nor within.  But all his sons [only] withstood evnom on the skin without.  Burhild brought another horn to Sinfittle and bade him drink, and everything passed as before.  And once more a third time she brought him the horn, and yet with accusatory words if he would not drink it up.  He spoke as before with Syemund; he said: “Let thy whisker-hairs sieve it, son!”  Sinfittle drank and was shortly dead.
Syemund bore him through long journeys in his arms and came to a single firth, narrow and long, and thereat there was one little ship and one man on it.  He offered Syemund passage around the firth.  But when Syemund bore the body out into the ship then the boat was fully loaded.  The churl said that Syemund should journey further inland along the firth.  The churl pushed out the ship and shortly disappeared.
King Syemund stayed for a long time in Denmark in Burhild’s realm after he got her hand.  Then Syemund journeyed south into Frankland to the realm he owned there.  Then he got the hand of Hardise, the daughter of King Eanlime.  Their son was Siward.  King Syemund fell in battle at the hands of the sons of Hunding.  But Hardise was then given to Atholf, the son of King Helpric.  There Siward grew up in his childhood.  Syemund and all his sons were far ahead of all other men in strength and size and thought and all accomplishments.  Siward was furthest of them all, and in ancient sources men call him ahead of all men and the most worshipful of raiding kings.\epb\epg


\bpg\bpa Grípir hét sonr Ęy·lima, bróðir Hjǫr·dísar.  Hann réð lǫndum ok vas allra manna vitrastr ok fram-víss.  Sig·urðr ręið ęinn saman ok kom til hallar Grípis.  Sig·urðr vas auð-kęnndr.  Hann hitti mann at máli úti fyr hǫll’inni; sá nęfndi⸗sk Gęitir.  Þá kvaddi Sig·urðr hann máls, ok spyrr:\epa

\bpb Griper was the name of Eanlime’s son, Hardise’s brother.  He ruled lands and was of all men the smartest, and forthwise.  Siward rode alone and came to Griper’s hall; Siward was easily recognized.  He went up to speak to a man outside of the hall; that one called himself Goater.  Then Siward greeted him with speech, and asks:\epb\epg

\subsection{Grípis spǫ́ (‘The Spae of Griper’)}

\bvg\bva%
\llap{„}Hvęrr \alst{b}yggir hér \hld\ \alst{b}orgir þessar? &
Hvat þann \edtrans{\alst{þ}jóð-konung}{great king}{\Bfootnote{An old poetic word, cognate with OE \emph{þéod-cyning} ‘id.’, OS \emph{þiod-kuning} ‘id.’
The first element is most likely to be equated with Gothic \emph{þiuþi-} ‘good’ (in \emph{þiuþi-qiss} ‘blessing’, calquing Greek \textgreek{εὐλογῐ́ᾱ}, \cite[s.v. \textit{þiuþiqiss}]{Streitberg1910}) rather than Gothic \emph{þiuda} ‘people, tribe’, ON \emph{þjóð} ‘id.’, OE \emph{þéod} ‘id.’

Its distribution in ON poetry is not altogether neutral.
In the Poetic Edda it appears, outside of the present poem, in \textlink{Sigurdskamma}[35]/3, 36/4, \textlink{Atlakvida}[47]/3, \textlink{Gudrunarhvot}[14]/2, and \textlink{Hamdismal}[4]/1, all likely old poems.
In Scaldic poetry it appears in the early \Ynglingatal\ 22 (ca. 850 \ce) and \Glymdrapa\ 7 (ca. 900 \ce), but appears to have been proscribed in the formal praise poetry of the 10th and earliest 11th century until it was revived by Sighwat Thurfrithson in the 1010s, from which point it is found in praise poetry until the extinction of the tradition in the 13th c.
In apocryphal Eddic poetry the word only appears in \emph{HjǪ 19} and Gunnleif \emph{Merl I} 58 (both in \Skp\ 8).}} \hld\ \alst{þ}egnar nęfna?“ &
\llap{„}\alst{G}rípir hęitir \hld\ \alst{g}umna stjóri, &
sá’s \alst{f}astri rę́ðr \hld\ \alst{f}oldu ok þegnum.“\eva

\bvb “{\huge W}\textsc{ho dwells here} in these strongholds? \\
What is this great king called by thanes?” \\
“Griper is the name of the steerer of men \\
who rules the steadfast land and thanes.”\evb\evg


\bvg\bva%
\llap{„}Es \alst{h}orskr konungr \hld\ \alst{h}ęima ï landi? &
Mun sá \alst{g}ramr við mik \hld\ \alst{g}anga at mę́la? &
\alst{M}áls es þarfi \hld\ \alst{m}aðr ȯ·kunnigr, &
vil’k \alst{f}ljót⸗liga \hld\ \alst{f}inna Grípi.“\eva

\bvb “Is the wise king at home in the land? \\
Will that chief go to speak with me? \\
In need of speech is this foreign man; \\
I wish shortly to find Griper.”\evb\evg


\bvg\bva%
\llap{„}Þęss mun \alst{g}laðr konungr \hld\ \alst{G}ęiti spyrja &
hvęrr sá \alst{m}aðr sé \hld\ es \alst{m}áls kveðr Grípi.“ &
\llap{„}\alst{S}ig·urðr ek heiti, \hld\ borinn \alst{S}ig·mundi &
ęn \alst{H}jǫr·dís es \hld\ \alst{h}ilmis móðir.“\eva

\bvb “The glad king will ask Goater \\
who the man is who asks Griper for speech.” \\
“Siward I am called, born to Syemund, \\
but Hardise is the helmet-bearer’s [my] mother.”\evb\evg


\bvg\bva%
Þȧ gekk \alst{G}ęitir \hld\ \alst{G}rípi at sęgja: &
\llap{„}Hér’s maðr \alst{ú}ti \hld\ \alst{ȯ}·kuðr kominn; &
hann ’s \alst{í}tar⸗ligr \hld\ at \alst{ȧ}·liti; &
sá vill, \alst{f}ylkir, \hld\ \alst{f}und þïnn hafa.“\eva

\bvb Then Goater went to tell Griper: \\
“Here outside a foreign man is come; \\
he is radiating of hue. \\
That one, O marshaller, wishes to have thy audience.”\evb\evg


\bvg\bva%
Gęngr ór \alst{sk}ála \hld\ \alst{sk}atna dróttinn &
ok \alst{h}ęilsar vel \hld\ \alst{h}ilmi komnum: &
\llap{„}Þigg þú hér, \alst{S}ig·urðr, \hld\ vę́ri \alst{s}ø̇mra fyrr! &
Ęn þú, \alst{G}ęitir, tak \hld\ við \edtrans{\alst{G}rana}{Grane}{\Bfootnote{Siward’s horse.}} sjǫlfum.“\eva

\bvb Out of his halls goes the lord of warriors \\
and well greets the helmet-bearer who has come. \\
“Here receive [welcome], Siward!  Sooner would have been more fit! \\
But thou, Goater, take care of Grane.”\evb\evg


\bvg\bva%
\alst{M}ę́la nǫ́mu \hld\ ok \alst{m}argt hjala &
þá’s \alst{r}áð-spakir \hld\ \alst{r}ekkar fundu⸗sk. &
\speakernote{[Sig·urðr:]}%
\llap{„}Sęg-ðu \alst{m}ér ef þú vęizt, \hld\ \alst{m}óður-bróðir, &
hvé mun \alst{S}ig·urði \hld\ \alst{s}núna ę́vi?“\eva

\bvb They took to speak and chatter much \\
when the council-wise champions found each other. \\
\speakernoteb{[Siward:]}%
“Tell me, if thou knowest, O mother’s brother: \\
how will Siward’s age turn out?”\evb\evg


\bvg\bva%
\llap{„}Þú \alst{m}unt \alst{m}aðr vesa \hld\ \alst{m}ę́tstr und sólu &
ok \alst{h}ę́str borinn \hld\ \alst{h}vęrjum jǫfri; &
\alst{g}jǫfull af \alst{g}olli \hld\ ęn \alst{g}løggr flugar, &
\alst{í}tr \alst{ȧ}·liti \hld\ ok ï \alst{o}rðum spakr.“\eva

\bvb \llap{„}Thou wilt be a man noblest beneath the sun, \\
and born higher than every ruler, \\
giving of gold but stingy of flight, \\
radiant of hue and wise in words.“\evb\evg


\bvg\bva%
\llap{„}Sęg, \alst{g}ęgn konungr, \hld\ \alst{g}ørr an ek spyrja &
\alst{s}notr \alst{S}ig·urði \hld\ ef þú \alst{s}já þikki⸗sk, &
hvat mun \alst{f}yrst gøra⸗sk \hld\ til \alst{f}arnaðar &
þȧ’s ór \alst{g}arði em’k \hld\ \alst{g}ęnginn þïnum?“\eva

\bvb “Tell it, O fine king, clearer than I might ask, \\
---O wise---to Siward, if thou seemest to see: \\
What will first happen to my advantage \\
when I am gone out of thy house?”\evb\evg


\bvg\bva%
\llap{„}\alst{F}yrst munt, \alst{f}ylkir, \hld\ \alst{f}ǫður of hęfna &
ok \edtrans{\alst{Ęy}·lima \hld\ alls harms}{all Eanlime’s injuries}{\Bfootnote{Eanlime was the father of Siward’s mother, Hardise, cf.~P1 above.}} reka; &
þú munt \alst{h}arða \hld\ \alst{H}undings sonu &
\alst{s}njalla fęlla--- \hld\ munt \alst{s}igr hafa.“\eva

\bvb “First wilt thou, marshaller, take revenge for thy father \\
and avenge all Eanlime’s injuries. \\
Harshly wilt thou fell the brave sons \\
of Hunding; thou wilt have victory.”\evb\evg


\bvg\bva%
\llap{„}Sęg, \alst{í}tr konungr, \hld\ \alst{ę́}ttingi mér &
\alst{h}ęldr \alst{h}orsk⸗liga \hld\ es vit \alst{h}ugat mę́lum; &
\alst{s}ér þú \alst{S}ig·urðar \hld\ \alst{s}nǫr brǫgð fyr\emph{ir} &
þau’s \alst{h}ę́st fara \hld\ und \alst{h}imin-skautum?“\eva

\bvb “Tell, O radiant king, to me, thy kinsman, \\
rather wisely, as we speak thoughtfully: \\
foreseest thou any of Siward’s nimble deeds \\
which fare highest under the corners of heaven?”\evb\evg


\bvg\bva%
\llap{„}Munt \alst{ęi}nn vega \hld\ \alst{o}rm inn frȧna &
þann’s \alst{g}rǫ́ðugr liggr \hld\ ȧ \alst{G}nita-hęiði; &
þú munt \alst{b}ǫ́ðum \hld\ at \alst{b}ana verða &
\alst{R}ęgin ok Fáfni; \hld\ \alst{r}étt sęgir Grípir.“\eva

\bvb “Thou alone wilt slay the gleaming Wyrm \\
which lies greedy on the Gnit-heath. \\
Thou wilt be the bane of them both: \\
Rein and Fathomer---Griper tells rightly.”\evb\evg


\bvg\bva%
\llap{„}\alst{Au}ðr mun \alst{ǿ}rinn \hld\ ef \alst{ę}fli’k svá &
\alst{v}íg með \alst{v}irðum \hld\ sem þú \alst{v}isst sęgir; &
\edtrans{\alst{l}ęið at·huga}{Give thy thought to it}{\Bfootnote{\emph{lęiða at·huga} ‘give one’s thought to something’ is a idiomatic \parenciteshorthand[s.v.~at·hugi]{ONP}.}} \hld\ ok \alst{l}ęngi sęg &
hvat mun \alst{ę}nn vesa \hld\ \alst{ę́}vi minnar?“\eva

\bvb “The wealth will be ample if I accomplish such \\
a killing amongst men as thou sayest for certain. \\
Give thy thought to it and speak for long: \\
what will further become of my age?”\evb\evg


\bvg\bva%
\llap{„}Þú munt \alst{f}inna \hld\ \alst{F}áfnis bǿli &
ok \alst{u}pp taka \hld\ \alst{au}ð inn fagra, &
\alst{g}olli hlǿða \hld\ ȧ \alst{G}rana bógu; &
ríðr til \alst{G}júka \hld\ \alst{g}ramr víg-risinn.“\eva

\bvb “Thou wilt find Fathomer’s dwelling \\
and take up the fair treasure, \\
load the gold on Grane’s shoulders; \\
thou wilt ride to Yivick a man proud of killing.”\evb\evg


\bvg\bva%
\llap{„}Ęnn skalt \alst{h}ilmi \hld\ ï \alst{h}ugaðs-rǿðu &
\alst{f}ram-lyndr jǫfurr \hld\ \alst{f}lęira sęgja; &
\alst{g}ęstr em’k \alst{G}júka \hld\ ok ek \alst{g}ęng þaðan, &
hvat mun \alst{ę}nn vesa \hld\ \alst{ę́}vi minnar?“\eva

\bvb “Further shalt thou to the marshaller in thoughtful conversation \\
O forth-minded chieftain, tell more things: \\
I am the guest of Yivick, and then I go thence; \\
what will further become of my age?”\evb\evg


\bvg\bva%
\llap{„}Søfr ȧ \alst{f}jalli \hld\ \alst{f}ylkis dóttir &
\alst{b}jǫrt ï \alst{b}rynju \hld\ ęptir \alst{b}ana Hęlga; &
þú munt \alst{h}ǫggva \hld\ \alst{h}vǫssu sverði, &
\alst{b}rynju rísta \hld\ með \edtrans{\alst{b}ana Fáfnis}{Fathomer’s bane}{\Bfootnote{Viz.~the sword Riddle; cf.~\textlink{Fafnismal}[P3].}}.“\eva

\bvb “On the fell sleeps the marshaller’s daughter, \\
bright in a byrnie, after Hallow’s killing. \\
Thou wilt cut it with a sharp sword, \\
carve up the byrnie with Fathomer’s bane.”\evb\evg


\bvg\bva%
\llap{„}\alst{B}rotin es \alst{b}rynja \hld\ \alst{b}rúðr mę́la tękr &
es \alst{v}aknaði \hld\ \alst{v}íf ór svefni; &
hvat mun \alst{s}nót at hęldr \hld\ við \alst{S}ig·urð mę́la &
þat ’s at \alst{f}arnaði \hld\ \alst{f}ylki verði?“\eva

\bvb “Broken is the byrnie; the bride takes to speak, \\
the woman who awakened from her sleep. \\
What will the wife thereafter speak unto Siward, \\
which will be to the marshaller’s advantage?”\evb\evg


\bvg\bva%
\llap{„}Hǫ̇n mun \alst{r}íkjum þér \hld\ \alst{r}u̇nar kęnna &
\alst{a}llar þę́r’s \alst{a}ldir \hld\ \alst{ęi}gna⸗sk skyldu &
ok ȧ \alst{m}anns tungu \hld\ \alst{m}ę́la hvęrja; &
\edtrans{\alst{l}if með \alst{l}ę́kning}{Live with leeching skill!}{\Bfootnote{I.e., medicinal skill, medicine.  Clearly alluding to \textlink{Sigrdrifumal}[4]/4: \emph{ok lę́knis-hęndr meðan lifum} ‘and a leecher’s hands, while we live’.}} \hld\ \alst{l}if hęill konungr!“\eva

\bvb “She will teach thee, O powerful one, runes, \\
all those which mankind should possess, \\
and on every man’s tongue be spoken. \\
Live with leeching skill! Live hale, O king!”\evb\evg


\bvg\bva%
\llap{„}\edtext{Þȧ’s því lokit \hld\ numin eru frǿði}{\Bfootnote{Alliteration is missing from this line and no obvious emendation can be found.}} &
ok em \alst{b}raut þaðan \hld\ \alst{b}úinn at ríða; &
\alst{l}ęið at·huga \hld\ ok \alst{l}ęngra sęg &
hvat mun \alst{m}ęirr vesa \hld\ \alst{m}innar ę́vi?“\eva

\bvb “Then that is come to an end; the learning has been taken, \\
and I am ready to ride away thence. \\
Give thy thought to it and speak for longer: \\
what will yet further become of my age?”\evb\evg


\bvg\bva%
\llap{„}Þú munt \alst{h}itta \hld\ \alst{H}ęimis bygðir &
ok \alst{g}laðr vera \hld\ \alst{g}ęstr þjóð-konungs; &
\alst{f}arit es, Sig·urðr, \hld\ þat’s ek \alst{f}yr·vissa’k, &
skal⸗a \alst{f}ręmr an svá \hld\ \alst{f}regna Grípi.“\eva

\bvb “Thou wilt find the farms of Homer \\
and be the glad guest of that great king. \\
Now that has passed, O Siward, which I foreknew; \\
further than so shalt thou not ask Griper.”\evb\evg


\bvg\bva%
\llap{„}Nú fǿr mér \alst{ę}kka \hld\ \alst{o}rð þat’s mę́ltir &
því’t langt \alst{f}ramm of sér \hld\ \alst{f}ylkir lęngra; &
vęitst \alst{o}f·mikit \hld\ \alst{a}ngr Sig·urði, &
því þú \alst{G}rípir þat \hld\ \alst{g}ørra sęgja!“\eva

\bvb “Now the word which thou hast spoken brings me grief; \\
for the marshaller foresees much further. \\
Thou knowest too great trouble for Siward; \\
therefore shalt thou, Griper, tell that more clearly!”\evb\evg


\bvg\bva%
\llap{„}Lá mer umb \alst{ø̇}sku \hld\ \alst{ę́}vi þinnar &
\alst{l}jósast fyr\emph{ir} \hld\ \alst{l}íta ęptir; &
\alst{r}étt em’k⸗a ek \hld\ \alst{r}áð-spakr taliðr &
ne in heldr \alst{f}ram-víss \hld\ \alst{f}arit þat’s ek vissa’k.“\eva

\bvb “Throughout the youth of thy age \\
it was brightest for me to behold. \\
I am not rightly counted counsel-wise, \\
nor indeed forthwise; passed is that which I have known.”\evb\evg


\bvg\bva%
\llap{„}\alst{M}ann vęit’k ęngi \hld\ fyr \alst{m}old ofan &
þann’s \alst{f}lęira sé \hld\ \alst{f}ramm an þú Grípir; &
skal⸗at \alst{l}ęyna \hld\ þó’tt \alst{l}jótt sér &
eða \alst{m}ęin gøri⸗sk \hld\ ȧ \alst{m}ïnum hag.“\eva

\bvb “I know no man above the earth \\
who foresees more things than thou, Griper. \\
Thou shalt not hide it, although thou seest ugliness \\
or evil befalling my state of life.”\evb\evg


\bvg\bva%
\llap{„}Es⸗a með \alst{l}ǫstum \hld\ \alst{l}ǫgð ę́vi þér; &
lát-tu, inn \alst{í}tri, \hld\ þat, \alst{ǫ}ðlingr, nema⸗sk, &
því at \alst{u}ppi mun \hld\ meðan \alst{ǫ}ld lifir, &
\alst{n}add-éls boði, \hld\ \alst{n}afn þitt vera.“\eva

\bvb “With blemishes is thy age not laid out; \\
let thyself, radiant athling, take that to heart, \\
for remembered will while mankind lives, \\
O beseecher of the sword-storm \ken{battle > warrior}, thy name be!”\evb\evg


\bvg\bva%
\llap{„}\alst{V}erst hyggjum því \hld\ \alst{v}erðr at skilja⸗sk &
\alst{S}ig·urðr við fylki \hld\ at \alst{s}ó-gǫru. &
\alst{l}ęið vísa þú \hld\ \alst{l}agt ’s allt fyr\emph{ir}, &
\alst{m}ę́rr, \alst{m}ér ef þú vilt \hld\ \alst{m}óður bróðir.“\eva

\bvb “We think it worst for Siward to part \\
from the marshaller after such a saying. \\
Show me the path---all is laid out ahead--- \\
if thou wilt, O famed mother’s brother!”\evb\evg


\bvg\bva%
\llap{„}Nú skal \alst{S}ig·urði \hld\ \alst{s}ęgja gǫrva &
alls \alst{þ}ęngill mik \hld\ til \alst{þ}ęss nęyðir; &
munt \alst{v}íst \alst{v}ita \hld\ at \alst{v}ę́t-ki lýgr, &
\alst{d}ǿgr ęitt es þér \hld\ \alst{d}auði ę́tlaðr.“\eva

\bvb “Now shall I tell Siward clearly \\
as the prince compels me to do it. \\
Thou shalt truly know that I nowise lie; \\
on a single half-day thy death is appointed.”\evb\evg


\bvg\bva%
\llap{„}Vil’k⸗at ek \alst{r}ęiði \hld\ \alst{r}íks þjóð-konungs &
\alst{g}óð-ráðs \edtext{at *}{\Afootnote{\emph{at at} \Regius}} hęldr \hld\ \alst{G}rípis þiggja; &
nú vill \alst{v}íst \alst{v}ita \hld\ þó’tt \alst{v}ilkit sé &
hvat á \alst{s}ẏnt \alst{S}ig·urðr \hld\ \alst{s}ér fyr hǫndum.“\eva

\bvb “TODO. \\
TODO. \\
TODO. \\
TODO”\evb\evg


\bvg\bva%
\llap{„}\alst{F}ljóð ’s at Hęimis \hld\ \alst{f}agrt ȧ·litum. &
hana \alst{B}ryn·hildi \hld\ \alst{b}ragnar nęfna; &
\alst{d}óttir Buðla \hld\ ęn \alst{d}ýrr konungr &
\alst{h}arð-úgðigt man \hld\ \alst{H}ęimir fǿðir.“\eva

\bvb “There is a lady at Homer’s, fair of hue; \\
her do the warriors call Byrnhild. \\
She is the daughter of Budle, but the esteemed king \\
Homer raises the hard-minded girl.”\evb\evg


\bvg\bva%
\llap{„}Hvat ’s \alst{m}ik at því \hld\ þó’tt \alst{m}ę́r sé &
\alst{f}ǫgr ȧ·liti \hld\ \alst{f}ǿdd at Hęimis? &
Þat skalt \alst{G}rípir \hld\ \alst{g}ǫrva sęgja &
því’t þú \alst{ǫ}ll of sér \hld\ \alst{ø}r·lǫg fyr\emph{ir}.“\eva

\bvb “What does it mean for me, though there be a maiden \\
fair of hue raised at Homer’s? \\
That shalt thou, Griper, tell clearly, \\
for thou foreseest all orlays.”\evb\evg


\bvg\bva%
\Ballnote{The description of Siward’s despair appears to be pulling from \textlink{Havamal}[114].}%
\llap{„}Hǫ̇n \alst{f}irrir þik \hld\ \alst{f}lęstu gamni &
\alst{f}ǫgr ȧ·liti \hld\ \alst{f}óstra Hęimis; &
\edtext{\alst{s}vefn þú né \alst{s}øfr \hld\ né umb \alst{s}akar dø̇mir}{\Bfootnote{The line closely parallels \textlink{GudrunTwo}[3]/3.}} &
gȧr⸗a þú \alst{m}anna \hld\ nema þú \alst{m}ęy sér.“\eva

\bvb “She will deprive thee of most pleasures, \\
the foster-daughter of Homer, fair of hue. \\
Thou wilt not sleep thy sleep nor preside over matters; \\
thou heedest no men if thou canst not see the maiden.”\evb\evg


\bvg\bva%
\llap{„}Hvat mun til \alst{l}íkna \hld\ \alst{l}agt Sig·urði &
\alst{s}ęg Grípir þat \hld\ ef \alst{s}já þikki⸗sk. &
\alst{m}un’k \alst{m}ęy ná \hld\ \alst{m}undi kaupa &
þȧ ina \alst{f}ǫgru \hld\ \alst{f}ylkis dóttur?“\eva

\bvb “What will be laid out to relive Siward; \\
tell, O Griper, that, if thou seemest to see it: \\
Will I get to buy the maiden for a bride-price, \\
that fair daughter of the marshaller?”\evb\evg


\bvg\bva%
\llap{„}\alst{I}t munuð \alst{a}lla \hld\ \alst{ęi}ða vinna &
\alst{f}ull \alst{f}ast⸗liga--- \hld\ \alst{f}á munuð halda; &
vesit hęfir þú \alst{G}júka \hld\ \alst{g}ęstr ęina nǫ́tt &
mant⸗at \alst{h}orska \hld\ \alst{H}ęimis fóstru.“\eva

\bvb “Ye two will swear all oaths, \\
fully, firmly; few will ye keep. \\
Thou wilt have been the guest of Yivick for a single night; \\
thou wilt not remember the wise foster-daughter of Homer.”\evb\evg


\bvg\bva%
\llap{„}Hvárt ’s þȧ \alst{G}rípir \hld\ \alst{g}ę̇tt þęss fyr mér &
sér þú \alst{g}ęð-lęysi \hld\ ï \alst{g}rams skapi &
es skal við \alst{m}ęy þá \hld\ \alst{m}ǫ́lum slíta &
es \alst{a}lls hugar \hld\ \alst{u}nna þȯttumk.“\eva

\bvb “TODO. \\
TODO. \\
TODO. \\
TODO”\evb\evg


\bvg\bva%
\llap{„}Þú verðr \alst{s}iklingr \hld\ fyr \alst{s}vikum annars &
munt \alst{G}rím·hildar \hld\ \alst{g}jalda ráða; &
mun \alst{b}jóða þér \hld\ \alst{b}jart-haddat man &
\alst{d}óttur sïna \hld\ \alst{d}ręgr hǫ̇n vel at gram.“\eva

\bvb “TODO. \\
TODO. \\
TODO. \\
TODO”\evb\evg


\bvg\bva%
\llap{„}Mun’k við þȧ \alst{G}unn·ar \hld\ \alst{g}ǫrva hlęyti &
ok \alst{G}uð·ru̇nu \hld\ \alst{g}anga at ęiga &
\alst{f}ull-kvę̇ni þȧ \hld\ \alst{f}ylkir vę́ri &
ef \alst{m}ęin-tregar \hld\ \alst{m}ér angraði-t.“\eva

\bvb “TODO. \\
TODO. \\
TODO. \\
TODO”\evb\evg


\bvg\bva%
\llap{„}Þik mun \alst{G}rím·hildr \hld\ \alst{g}ǫrva véla &
mun hǫ̇n \alst{B}ryn·hildar \hld\ \alst{b}iðja fẏsa &
\alst{G}unn·ari til handa \hld\ \alst{G}otna dróttni &
hęitr þú \alst{f}ljót⸗liga \alst{f}ǫr \hld\ \alst{f}ylkis móður.“\eva

\bvb “TODO. \\
TODO. \\
TODO. \\
TODO”\evb\evg


\bvg\bva%
\llap{„}\alst{M}ęin eru fyr hǫndum \hld\ \alst{m}á’k líta þat &
\alst{r}atar gǫr⸗liga \hld\ \alst{r}áð Sig·urðar, &
ef ek skal \alst{m}ę́rrar \hld\ \alst{m}ęyjar biðja &
\alst{ǫ}ðrum til handa \hld\ þęirar ek \alst{u}nna vel.“\eva

\bvb “TODO. \\
TODO. \\
TODO. \\
TODO”\evb\evg


\bvg\bva%
\llap{„}\alst{É}r munuð \alst{a}llir \hld\ \alst{ęi}ða vinna &
\alst{G}unn·arr ok Hǫgni \hld\ ęn þú \alst{g}ramr þriði, &
því’t \alst{l}itum víxla \hld\ es ȧ \alst{l}ęið eruð &
\alst{G}unn·arr ok þú--- \hld\ \alst{G}rípir lýgr ęigi.“\eva

\bvb “TODO. \\
TODO. \\
TODO. \\
thou and Guther---Griper lies not.”\evb\evg


\bvg\bva%
\llap{„}\alst{H}ví gęgnir þat \hld\ \alst{h}ví skulum skipta &
\alst{l}itum ok \alst{l}ǫ́tum \hld\ es ȧ \alst{l}ęið erum. &
þar mun \alst{f}lá-rǿði \hld\ \alst{f}ylgja annat &
\alst{a}talt með \alst{ǫ}llu--- \hld\ \alst{ę}nn sęg Grípir!“\eva

\bvb “TODO. \\
TODO. \\
TODO. \\
TODO”\evb\evg


\bvg\bva%
\llap{„}\alst{L}it hęfir þú Gunn·ars \hld\ ok \alst{l}ę́ti hans, &
\alst{m}ę́lsku þïna \hld\ ok \alst{m}ęgin-hyggjur; &
munt \alst{f}astna þér \hld\ \alst{f}ram-lundaða &
\alst{f}óstru Hęimis \hld\ sér vę́tr \alst{f}yr því.“\eva

\bvb “Thou hast Guther’s hue and his voice, \\
thy own speech and intelligence; \\
TODO. \\
TODO”\evb\evg


\bvg\bva%
\llap{„}\alst{V}erst hyggjum því \hld\ \alst{v}ándr mun’k hęitinn &
\alst{S}ig·urðr með \alst{s}ęggjum \hld\ at \alst{s}ó-gǫru; &
\alst{v}ilda’k ęigi \hld\ \alst{v}élum bęita &
\alst{jǫ}fra brúðr \hld\ es \alst{ǿ}ðsta vęit’k.“\eva

\bvb “TODO. \\
TODO. \\
TODO. \\
TODO”\evb\evg


\bvg\bva%
\llap{„}Þú munt \alst{h}víla, \hld\ \alst{h}ęrs odd-viti, &
\alst{m}ę́rr hjá \alst{m}ęyju \hld\ sem þïn \alst{m}óðir sé; &
því mun \alst{u}ppi \hld\ meðan \alst{ǫ}ld lifir, &
\alst{þ}jóðar \alst{þ}ęngill, \hld\ \alst{þ}itt nafn vera.\eva

\bvb “Thou wilt rest, O famous point-knower of the host \ken{warrior}, \\
beside the maiden as if she were thy mother. \\
For that will remembered while mankind lives, \\
O prince of the nation, thy name be.\evb\evg


\bvg\bva%
Saman munu \alst{b}ru·llaup \hld\ \alst{b}ę́ði drukkin &
\alst{S}ig·urðar ok Gunn·ars \hld\ ï \alst{s}ǫlum Gjúka; &
þȧ \alst{h}ǫmum víxlið \hld\ es it \alst{h}ęim komið, &
\alst{h}ęfir \alst{h}vęrr fyr því \hld\ \alst{h}yggju sïna.“\eva

\bvb “Both weddings will be drunk together \\
---Siward’s and Guther’s---in the halls of Yivick. \\
Then ye will shift your hames when ye two come home; \\
TODO.”\evb\evg


\bvg\bva%
\llap{„}Mun \alst{g}óða kvǫ́n \hld\ \alst{G}unn·arr ęiga &
\alst{m}ę́rr með \alst{m}ǫnnum, \hld\ ---\alst{m}ér sęg Grípir!--- &
þó’tt hafi \alst{þ}rjár nę́tr \hld\ \alst{þ}egns brúðr hjá mér &
\alst{s}nar-lynd \alst{s}ofit? \hld\ \alst{S}·líks eru-t dø̇mi!\eva

\bvb “Will Guther have that good woman, \\
the renowned man amongst men---tell me Griper!--- \\
even though she has slept by me for three nights, \\
the sharp-minded thane’s bride?  Such a thing has no precedents!”\evb\evg


\bvg\bva%
Hvé mun at \alst{y}nði \hld\ \alst{ę}ptir verða &
\alst{m}ę́gð með \alst{m}ǫnnum? \hld\ \alst{M}ér sęg Grípir! &
Mun \alst{G}unn·ari \hld\ til \alst{g}amans ráðit &
\alst{s}íðan verða \hld\ eða \alst{s}jǫlfum mér?“\eva

\bvb “TODO. \\
TODO. \\
TODO. \\
TODO”\evb\evg


\bvg\bva%
\llap{„}\alst{M}innir þik ęiða \hld\ \alst{m}átt þęgja þó &
ann \alst{G}uð·ru̇nu \hld\ \alst{g}óðra ráða &
ęn \alst{B}ryn·hildr þikki⸗sk \hld\ \alst{b}rúðr vas gefin &
\alst{s}nót fiðr vélar \hld\ \alst{s}ér at hęfndum.“\eva

\bvb “TODO. \\
TODO. \\
TODO. \\
TODO”\evb\evg


\bvg\bva%
\llap{„}Hvat mun at \alst{b}ótum \hld\ \alst{b}rúðr sú taka &
es \alst{v}élar \alst{v}ér \hld\ \alst{v}ífi gerðum; &
hęfir \alst{s}nót af mér \hld\ \alst{s}varna ęiða &
\alst{ø}nga \alst{ę}fnda \hld\ en \alst{u}nat lítit?“\eva

\bvb “TODO. \\
TODO. \\
TODO. \\
TODO”\evb\evg


\bvg\bva%
\llap{„}Mun hǫ̇n \alst{G}unn·ari \hld\ \alst{g}ǫrva sęgja &
at þú \alst{ęi}gi vel \hld\ \alst{ęi}ðum þyrmir &
þȧ’s \alst{í}tr konungr \hld\ af \alst{ǫ}llum hug &
\alst{G}júka arfi \hld\ ȧ \alst{g}ram trúði.“\eva

\bvb “To Guther will she clearly tell \\
that thou dost not well respect the oaths, \\
while that radiant king, out of his whole heart, \\
---the heir of Yivick---trusted the chief [thee].”\evb\evg


\bvg\bva%
\llap{„}Hvat ’s þȧ, \alst{G}rípir? \hld\ \alst{G}et þęss fyr mér; &
mun’k \alst{s}aðr vesa \hld\ at \alst{s}ǫgu þęiri, &
eða \alst{l}ýgr ȧ mik \hld\ \alst{l}of-sę́l kona &
ok ȧ \alst{s}jalfa sik? \hld\ \alst{S}ęg Grípir þat!“\eva

\bvb “What happens then, Griper?  Reveal it to me; \\
will I be truthfully accused according to that speech, \\
or does the woman blessed of praise lie about me \\
and about her very self?  Tell, O Griper, this!”\evb\evg


\bvg\bva%
\llap{„}Mun fyr \alst{r}ęiði \hld\ \alst{r}ík brúðr við þik &
né af \alst{o}f·trega \hld\ \alst{a}ll-vel skipa. &
viðr þú \alst{g}óðri \hld\ \alst{g}rand aldri⸗gi &
þó ’s \alst{v}íf konungs \hld\ \alst{v}élum bęittuð.“\eva

\bvb “TODO. \\
TODO. \\
TODO. \\
TODO”\evb\evg


\bvg\bva%
\llap{„}Mun \alst{h}orskr Gunn·arr \hld\ at \alst{h}vǫtun hęnnar &
\alst{G}uð·þormr ok Hǫgni \hld\ \alst{g}anga síðan? &
Munu \alst{s}ynir Gjúka \hld\ af \alst{s}ifjungum mér &
\alst{ę}ggjar rjóða? \hld\ \alst{Ę}nn sęg Grípir!“\eva

\bvb “Will the wise Guther at her instigation, \\
with Godthorm and Hawn thereafter go? \\
Will the sons of Yivick, to break their kinship with me, \\
redden their blades?  Tell further, Griper!”\evb\evg


\bvg\bva%
\llap{„}Þȧ ’s \alst{G}uð·ru̇nu \hld\ \alst{g}rimmt umb hjarta &
\alst{b}rǿðr hęnnar þér \hld\ til \alst{b}ana ráða, &
ok at \alst{ø}ngu verðr \hld\ \alst{y}nði síðan &
\alst{v}itru \alst{v}ífi, \hld\ \alst{v}ęldr því Grím·hildr.\eva

\bvb “TODO. \\
TODO. \\
TODO. \\
TODO”\evb\evg


\bvg\bva%
Því skal \alst{h}ugga þik, \hld\ \alst{h}ęrs odd-viti, &
sú mun \alst{g}ipt lagit \hld\ ȧ \alst{g}rams ę́vi; &
mun⸗at \alst{m}ę́tri \alst{m}aðr \hld\ ȧ \alst{m}old koma &
und \alst{s}ólar \alst{s}jǫt \hld\ an, \alst{S}ig·urðr, þikkir.“\eva

\bvb For that shall she soothe(?) thee, O point-knower of the host; \\%TODO: "soothe"??
she will have laid venom in the ruler’s age. \\
No nobler man will come upon the earth \\
beneath the seat of the sun \ken{sky/heaven}, than thou, Siward, dost seem!\evb\evg


\bvg\bva%
\llap{„}\alst{Sk}iljum⸗k hęilir; \hld\ mun⸗at \alst{sk}ǫpum vinna! &
Nú hęfir þú, \alst{G}rípir, vęl \hld\ \alst{g}ørt sem bęidda’k; &
\alst{f}ljótt myndir þú \hld\ \alst{f}ríðri sęgja &
\alst{m}ïna ę́vi \hld\ ef þú \alst{m}ę́ttir þat!“\eva

\bvb “Let us part hale; man will not withstand the \inx[C]{shape}[shapes]! \\
Now hast thou, Griper, well done as I asked. \\
Sooner wouldst thou have told finer things \\
of my age, if thou couldst do so!”\evb\evg

\sectionline
